\PassOptionsToPackage{table,xcdraw}{xcolor}
\documentclass[11pt, a4paper]{article}

\usepackage[T1]{fontenc}
\usepackage[utf8]{inputenc}
\usepackage{tgpagella}
\usepackage[spanish, es-noshorthands]{babel}
\usepackage{amsmath, amssymb}
\usepackage{booktabs}
\usepackage{tabularx}
\usepackage{multicol}
\usepackage[most]{tcolorbox}
\usepackage[american]{circuitikz}
\usepackage[margin=2cm]{geometry}

\definecolor{ucbblue}{RGB}{0, 51, 102}

\begin{document}

\begin{center}
    \Large\textbf{\textcolor{ucbblue}{Suplemento Lab 4: Arquitecturas de Alimentación y Verificación Experimental}}[cite: 14]
\end{center}

\begin{tcolorbox}[colback=ucbblue!5, colframe=ucbblue, title=\textbf{Objetivo de la Sesión[cite: 14]}]
    No estandarizaremos la fuente física; estandarizaremos el marco eléctrico y la evidencia experimental[cite: 14]. Su grupo debe reportar qué alimentación usa, qué referencia rige el circuito, qué rieles efectivos ve el Op-Amp y cómo afecta esto a sus mediciones dinámicas[cite: 14].
\end{tcolorbox}

\section*{1. Arquitecturas de Alimentación Válidas[cite: 14]}

Existen tres arquitecturas físicas válidas para este laboratorio. \textbf{El estudiante debe medir y documentar los rieles reales[cite: 14].}

\vspace{0.3cm}
\noindent\textbf{Arquitectura A: Fuente Dual de Banco[cite: 14]}
\begin{itemize}
    \setlength{\itemsep}{0pt}
    \item \textbf{Referencia:} $0\,\text{V}$[cite: 14].
    \item \textbf{Rieles:} $V_+ > 0$ y $V_- < 0$ (Nominalmente $\pm 9\,\text{V}$ a $\pm 10\,\text{V}$)[cite: 14].
\end{itemize}

\vspace{0.3cm}
\noindent\textbf{Arquitectura B: Dos Baterías de 9 V[cite: 14]}

\begin{minipage}[t]{0.48\textwidth}
    \vspace{0pt}
    Utilice dos baterías independientes de $9\,\text{V}$. El punto de unión entre ambas forma la referencia del circuito ($0\,\text{V}$).
    
    \vspace{0.3cm}
    \textbf{Conexión:}
    \begin{itemize}
        \item Batería 1 (+): Conectar a riel $V_+$.
        \item Batería 1 (-): Conectar a la unión central ($0\,\text{V}$).
        \item Batería 2 (+): Conectar a la unión central ($0\,\text{V}$).
        \item Batería 2 (-): Conectar a riel $V_-$.
    \end{itemize}
    \begin{tcolorbox}[colback=red!5, colframe=red!70!black]
        \textbf{Nota Crítica:} El punto medio es la referencia $0\,\text{V}$. \textbf{NO} trate el terminal negativo de la batería inferior como "tierra".
    \end{tcolorbox}
\end{minipage}
\hfill
\begin{minipage}[t]{0.48\textwidth}
    \vspace{0pt}
    \begin{center}
        \begin{circuitikz}[scale=0.9, transform shape, thick]
            % Bateria 1 (Superior)
            \draw (0,3) to[battery2, l=$9\,\text{V}$ (Bat 1)] (0,1);
            \draw (0,3) to[short, *-o] (2,3) node[right]{$V_+ \approx +9\,\text{V}$};
            
            % Union central y tierra
            \draw (0,1) to[short, *-o] (2,1) node[right]{$0\,\text{V}$ (Midpoint)};
            \draw (0,1) node[ground]{};
            
            % Bateria 2 (Inferior)
            \draw (0,1) to[battery2, l=$9\,\text{V}$ (Bat 2)] (0,-1);
            \draw (0,-1) to[short, *-o] (2,-1) node[right]{$V_- \approx -9\,\text{V}$};
        \end{circuitikz}
    \end{center}
\end{minipage}

\vspace{0.3cm}
\noindent\textbf{Arquitectura C: Fuente Simple + Referencia Virtual Bufferizada[cite: 14]}

\begin{center}
    \begin{circuitikz}[scale=0.9, transform shape, thick]
        \draw (0,3) node[vcc]{$V_S \approx 18\,\text{V}$} 
              to[R, l={$R_1$}] (0,1.5) coordinate(vm) 
              to[R, l={$R_2$}] (0,0) node[ground]{$0\,\text{V}$};
        \draw (vm) to[short, *-] ++(1,0) node[op amp, noinv input up, anchor=+] (buffer) {};
        \draw (buffer.-) -- ++(0,-1) coordinate(fbb) -| (buffer.out);
        \draw (buffer.out) to[short, *-o] ++(1.5,0) node[right]{$V_{REF}$};
        \draw (buffer.up) -- ++(0,0.3) node[above]{$V_S$};
        \draw (buffer.down) -- ++(0,-0.3) node[below]{$0\,\text{V}$};
        \node[below, align=center] at (1.5,-1) {\textit{Buffer LM358N}};
        \node[left] at (0,1.5) {$V_M$};
    \end{circuitikz}
\end{center}

\textbf{Derivación del Modelo[cite: 14]:} \\
A partir de una fuente física $V_S \approx 18\,\text{V}$[cite: 14], un divisor de tensión ideal con $R_1 = R_2$ produce un voltaje en el punto medio $V_M$[cite: 14]:
\begin{align*}
    V_M &= V_S \frac{R_2}{R_1 + R_2} \text{[cite: 14]} \\
    V_M &= V_S \frac{R_2}{2R_2} = \frac{V_S}{2} \text{[cite: 14]}
\end{align*}
La resistencia de Thévenin del divisor es:
\begin{equation*}
    R_{TH} = R_1 \parallel R_2 = \frac{R_1 R_2}{R_1 + R_2} \text{[cite: 14]}
\end{equation*}
Si el punto medio alimentara el circuito directamente, cualquier corriente $I_{REF}$ extraída desplazaría el voltaje en $\Delta V \approx I_{REF} R_{TH}$[cite: 14]. Para evitar esto, se utiliza un seguidor de voltaje (buffer) con un LM358N[cite: 14]:
\begin{equation*}
    V_{REF} \approx V_M \approx \frac{V_S}{2} \text{[cite: 14]}
\end{equation*}

\textbf{Rieles Efectivos para la Arquitectura C[cite: 14]:} \\
El circuito principal toma a $V_{REF}$ como su "tierra virtual" ($0\,\text{V}_{circuito}$)[cite: 14]. Por lo tanto, los rieles efectivos que ve el amplificador operacional bajo prueba son[cite: 14]:
\begin{align*}
    V_{+,eff} &= V_S - V_{REF} \text{[cite: 14]} \\
    V_{-,eff} &= 0 - V_{REF} = -V_{REF} \text{[cite: 14]}
\end{align*}
\textit{Ejemplo Numérico:} Si $V_S = 17.8\,\text{V}$ y se mide $V_{REF} = 8.85\,\text{V}$[cite: 14]:
\begin{align*}
    V_{+,eff} &= 17.8\,\text{V} - 8.85\,\text{V} = \boxed{+8.95\,\text{V}} \text{[cite: 14]} \\
    V_{-,eff} &= 0\,\text{V} - 8.85\,\text{V} = \boxed{-8.85\,\text{V}} \text{[cite: 14]}
\end{align*}
\textbf{Nota de Estabilidad:} Mida $V_{REF}$ en vacío ($V_{REF,unloaded}$) y con el circuito conectado ($V_{REF,loaded}$)[cite: 14]. Si hay un cambio significativo ($\Delta V_{REF}$), alterará el punto de operación dinámico[cite: 14].

\section*{2. Riesgo Crítico de Conexión a Masa (Instrument Ground)[cite: 14]}
\begin{tcolorbox}[colback=red!5, colframe=red!70!black, title=\textbf{Verificación Obligatoria de Equipos}]
    \raggedright
    Los terminales de masa ("tierra") del osciloscopio o generador de funciones suelen estar conectados internamente a la tierra de protección de la red eléctrica[cite: 14]. 
    Si conecta esta masa de osciloscopio a $V_{REF}$ en la Arquitectura C, podría \textbf{cortocircuitar} la referencia virtual hacia tierra física, invalidando el experimento y dañando equipos[cite: 14]. \\
    \textbf{Acción:} Solicite aprobación docente \textbf{antes} de conectar las pinzas de masa si usa la Arquitectura C[cite: 14].
\end{tcolorbox}

\section*{3. Validación y Diagnóstico (Troubleshooting)[cite: 14]}
Si observa comportamiento anómalo en sus experimentos, siga este flujo en orden estricto[cite: 14]:
\begin{sloppypar}
\begin{enumerate}
    \setlength{\itemsep}{0pt}
    \item \textbf{Verificar entrada y swing de salida:} ¿$V_{in,pp}$ es realmente lo asumido?[cite: 14] ¿La onda está chocando con los rieles efectivos (clipping)?[cite: 14]
    \item \textbf{Verificar estabilidad de la referencia:} Para Arquitectura C, mida si $V_{REF}$ se ha movido[cite: 14].
    \item \textbf{Predecir GBW:} Compare $f$ con $f_{BW,pred} = GBW/A_N$[cite: 14]. Evidencia de limitación por GBW: atenuación suave y desfase[cite: 14].
    \item \textbf{Calcular demanda SR:} $SR_{req} = 2\pi f V_p$[cite: 14]. Evidencia: triangularización de la onda[cite: 14].
    \item \textbf{Revisar Instrumentación:} ¿Está el generador u osciloscopio atenuando la señal por sus propios límites?[cite: 14]
    \item \textbf{Atribuir causa:} Solo después de descartar lo anterior, atribuya el problema a GBW, SR o Clipping[cite: 14].
\end{enumerate}
\end{sloppypar}

\end{document}
